% e-TeX's `\middle`. §1191 in the e-TeX change file:
%
%     p:=new_noad; type(p):=t; scan_delimiter(delimiter(p),false);
%     if t=middle_noad then
%       begin type(p):=right_noad; subtype(p):=middle_noad;
%       end;
%     if t=left_noad then q:=p
%     else  begin q:=fin_mlist(p); unsave; end;
%     if t<>right_noad then
%       begin push_math(math_left_group); link(head):=q; tail:=p;
%       delim_ptr:=p;
%       end
%
% `\middle` is a `right_noad` carrying the subtype `middle_noad`. It closes the
% `math_left_group` its `\left` opened and immediately opens another whose list
% *is* the old one with the new noad appended — so the mlist that reaches §762
% is flat:
%
%     left_noad … middle … right_noad
%
% and §762 measures the whole of it, growing all three delimiters to the same
% `max_h`/`max_d`. That is the point: `\middle|` is as tall as the fences around
% it, where a bare `|` is one glyph of one size.
%
% Undefined here until now, so `\left\langle s_i \middle| s_i^2=1
% \right\rangle` raised "Undefined control sequence" and set the `|` as an
% ordinary character. 2608.18024v1 writes exactly that.
%
% This case names an e-TeX primitive, so `tools/showbox_check.sh` will not find
% it under `tex` and falls back to `pdftex -ini -etex` — which is the only
% engine that can answer, `\middle` not being TeX82 at all.
\catcode`\{=1 \catcode`\}=2 \catcode`\$=3 \catcode`\^=7 \catcode`\_=8
\font\r=cmr10 \r \hsize=300pt \parindent=0pt
\font\mi=cmmi10 \font\sy=cmsy10 \font\ex=cmex10
\textfont0=\r \scriptfont0=\r \scriptscriptfont0=\r
\textfont1=\mi \scriptfont1=\mi \scriptscriptfont1=\mi
\textfont2=\sy \scriptfont2=\sy \scriptscriptfont2=\sy
\textfont3=\ex \scriptfont3=\ex \scriptscriptfont3=\ex
\showboxdepth=6 \showboxbreadth=20 \tracingonline=0

% The shape from the paper: a tall-ish body, so the fences and the middle bar
% all have to grow past one glyph.
\setbox0=\hbox{$\left\langle a^2 \middle| b_j \right\rangle$}
\showbox0

% Nothing tall inside: every delimiter stays at its smallest size, and the
% middle one must not grow on its own.
\setbox0=\hbox{$\left( x \middle| y \right)$}
\showbox0

% Two of them in one pair.
\setbox0=\hbox{$\left( x \middle| y \middle| z \right)$}
\showbox0

% A `\middle` whose delimiter is a null one — `\middle.` — which has no glyph
% at all but still takes part in the sizing.
\setbox0=\hbox{$\left( x \middle. y \right)$}
\showbox0

% Grown by a fraction rather than by a script, so `max_h`/`max_d` come from a
% box and not from a character.
\setbox0=\hbox{$\left[ {a \over b} \middle| c \right]$}
\showbox0
\end
